\documentclass[10pt,a4paper]{amsart}
\usepackage[margin=19mm]{geometry}
\usepackage{amsmath,amssymb,mathtools}
\usepackage[scr=rsfs,cal=euler]{mathalfa}
\usepackage{xfrac}
\usepackage[unicode,hidelinks,bookmarks=false]{hyperref}
\newcommand{\ie}{i.e.\ }
\newcommand{\cf}{cf.\ }
\newcommand{\resp}{respectively\ }
\newcommand{\Subsection}{\subsection}
\providecommand{\nobreakdash}{\nobreak\hbox{-}}
\setlength{\emergencystretch}{2em}
\multlinegap0pt
\usepackage{bm}
\usepackage{bbm}
\usepackage{dsfont}
\usepackage{xspace}
\usepackage{cancel}
\usepackage{mathtools}
\usepackage{amsthm}
\makeatletter
\@ifundefined{thm}{\newtheorem{thm}{Thm}}{}
\@ifundefined{lem}{\newtheorem{lem}[thm]{Lemma}}{}
\makeatother
\newcommand{\E}{\mathds{E}}
\newcommand{\dP}{\mathbb{P}}
\newcommand{\dR}{\mathbb{R}}
\title[A Gaussian process limit for the self-normalized Ewens-Pitma]{A Gaussian process limit for the self-normalized Ewens-Pitman process}
\author{Bernard Bercu, Stefano Favaro}
\date{Source: arXiv:2601.11216v1 (CC BY 4.0)}
\begin{document}
\maketitle

\noindent\emph{Source and licence.} A Gaussian process limit for the self-normalized Ewens-Pitman process. Version arXiv:2601.11216v1 (CC BY 4.0), distributed under the Creative Commons Attribution 4.0 International licence, as recorded in the arXiv licence metadata for this exact version and shown on the abstract page https://arxiv.org/abs/2601.11216v1. Full rights notice: licenses/arxiv-2601.11216-CC-BY-4.0.txt.\par\smallskip

\noindent\textit{Editorial scope.} Counted unit: the Lem whose source label is \texttt{L-KEY} in the article (source0.tex lines 452--459, source label \texttt{L-KEY}), delivered together with the article's own complete original proof of it (source0.tex lines 461--605, one \texttt{proof} environment of 145 lines). No mathematics is written, repaired or re-derived: statement and proof are the article's own text line for line. Required context retained uncounted: eq:kalpha (lines 299--301); eq:kralpha (lines 303--305); DEFK (lines 403--405); DEFPN (lines 407--410); DEFKr (lines 412--414); DEFXIKr (lines 416--423); DEFIPr (lines 425--428); DEFAr (lines 432--435); DEFMr (lines 437--440); DEFSNr (lines 442--445); DEFbr (lines 447--450). Every definition, display and label of those lines is retained unchanged. Omitted from this package and not redistributed: 1--298, 302--302, 306--402, 406--406, 411--411, 415--415, 424--424, 429--431, 436--436, 441--441, 446--446, 451--451, 460--460, 606--1822. Nothing inside a retained excerpt is omitted or altered: every retained body is delivered exactly as the article's lines are written, so no omission receipt of any kind accompanies this package. Citations: the retained excerpts carry no citation command at all, so no bibliography entry of the article is transcribed and no reference is redistributed. Reference adaptation: none was needed and none was made; every reference inside the delivered unit resolves inside it, so no unresolved reference or citation is produced. Printed slips kept literal: no printed word, symbol, hypothesis or number of the counted statement or of its proof was altered. Layout and class: the article's own class is \texttt{article} with packages including geometry, amssymb, amsmath, amsthm, mathtools, tasks, bm, lscape, bbm, mathrsfs, dsfont, wasysym, hyperref, natbib; the wrapper is the lane's pinned \texttt{amsart} editorial wrapper at 10pt on A4, which restates the theorem-like environments the retained text needs and adds the article's own macro definitions that the retained text needs (\texttt{\textbackslash E}, \texttt{\textbackslash dP}, \texttt{\textbackslash dR}), with extra wrapper packages (bm, bbm, dsfont, xspace, cancel, mathtools). The delivered numbering is the unit's own. Assets: no retained span contains a figure environment, a graphics-inclusion command, a PDF-page inclusion, a table or a TikZ picture; no image file is copied into this package, the package declares no asset dependency and no image file is redistributed with it. Licence: version arXiv:2601.11216v1 of this article is distributed under the Creative Commons Attribution 4.0 International licence, as recorded in the arXiv licence metadata for this exact version and shown on the abstract page https://arxiv.org/abs/2601.11216v1. A full rights notice is delivered at licenses/arxiv-2601.11216-CC-BY-4.0.txt. Characters: every retained excerpt is pure ASCII, so the delivered engine, XeLaTeX with its default Latin Modern text font, reports no missing character.
\begin{equation}\label{eq:kalpha}
\lim_{n\rightarrow+\infty}\frac{K_{n}}{n^{\alpha}}=S_{\alpha,\theta}\qquad\text{a.s.},
\end{equation}

\begin{equation}\label{eq:kralpha}
\lim_{n\rightarrow+\infty}\frac{K_{r,n}}{n^{\alpha}}=p_{\alpha}(r) S_{\alpha,\theta}\qquad\text{a.s.},
\end{equation}

\begin{equation}\label{DEFK}
K_{n+1} = K_n + \xi_{n+1},
\end{equation}

\begin{equation}
\label{DEFPN}
p_n=\frac{\alpha K_n+ \theta}{n+\theta}.
\end{equation} 

\begin{equation}\label{DEFKr}
K_{r,n+1}=K_{r,n}+\xi_{r,n+1},
\end{equation}

\begin{equation}
\label{DEFXIKr}
   \dP(\xi_{r,n+1}=k\,|\,\mathcal{F}_{r,n})=\left \{ \begin{array}{ccc}
    {p_{r,n}}  & \text{ if } & \ \ k=1 \vspace{1ex}\\
    {q_{r,n}}  & \text{ if } & \ \ k=-1 \vspace{1ex}\\
     {1-p_{r,n}-q_{r,n}} & \text{ if }  & \ \ k=0,
   \end{array}  \right.
\end{equation}

\begin{equation}
\label{DEFIPr}
p_{r,n}=\frac{(r-1-\alpha)K_{r-1,n}}{n+\theta}\qquad \text{and}\qquad q_{r,n}=\frac{(r-\alpha)K_{r,n}}{n+\theta}.
\end{equation}

\begin{equation}
\label{DEFAr}
a_{r,n}=\prod_{k=r}^{n-1}\left(\frac{k+\theta}{k+\alpha+\theta-r}\right)=\frac{(\theta+1)^{(n-1)}}{(\alpha+\theta-r+1)^{(n-1)}}.
\end{equation}

\begin{equation}
\label{DEFMr}
M_{r,n}=a_{r,n}\left(S_{r,n}+ \frac{(-1)^{r} p_{\alpha}(r) \theta}{\alpha}\right)
\end{equation}

\begin{equation}
\label{DEFSNr}
S_{r,n}= \sum_{i=1}^{r}b_{r,i}K_{i,n}+(-1)^r p_{\alpha}(r) K_{n},
\end{equation}

\begin{equation}
\label{DEFbr}
b_{r,i}=(-1)^{r-i}\frac{(i-\alpha)^{(r-i)}}{(r-i)!}.
\end{equation}

\begin{lem}
\label{L-KEY}
For all $r \geq 1$, $(M_{r,n})$ is a locally square integrable martingale with predictable quadratic variation $\langle M_r \rangle_{n}$ satisfying
\begin{equation}
\label{LIMIPINMr}
\lim_{n\rightarrow+\infty}\frac{\langle M_{r}\rangle_{n}}{n^{2r-\alpha}}=p_{\alpha}(r)\frac{(r-\alpha)^{(r)}}{r!}\left(\frac{\Gamma(\alpha+\theta-r+1)}{\Gamma(\theta+1)}\right)^{2}S_{\alpha,\theta}\qquad\text{a.s.}
\end{equation}
\end{lem}

\begin{proof}
It follows from \eqref{DEFK}, \eqref{DEFKr}, \eqref{DEFXIKr} and \eqref{DEFSNr} that
\begin{align}
\label{eq:esr0}
\E[S_{r,n+1}\,|\,\mathcal{F}_{r,n}]
&=\E\left[\sum_{i=1}^{r}b_{r,i}K_{i,n+1}+(-1)^r p_{\alpha}(r) K_{n+1}\,\Bigr|\,\mathcal{F}_{r,n}\right] \qquad \text{a.s.}
\notag\\
&=\E\left[\sum_{i=1}^{r}b_{r,i}(K_{i,n}+\xi_{i,n+1})+(-1)^r p_{\alpha}(r)(K_{n}+\xi_{n+1})\,\Bigr|\,\mathcal{F}_{r,n}\right]  \qquad \text{a.s.} \notag \\
&=S_{r,n}+ \sum_{i=1}^{r}b_{r,i}(p_{i,n}-q_{i,n}) + (-1)^rp_{\alpha}(r) p_{n}  \qquad \text{a.s.}
\end{align}
Consequently, as $p_{1,n}=p_n$ and for all $i \geq 2$, $p_{i,n}=q_{i-1,n}$, we obtain from \eqref{eq:esr0} that
\begin{equation}
\label{eq:esr1}
\E[S_{r,n+1}\,|\,\mathcal{F}_{r,n}]=S_{r,n}+((-1)^rp_{\alpha}(r)+b_{r,1})p_{n}-q_{r,n}+\sum_{i=1}^{r-1}(b_{r,i+1}-b_{r,i})q_{i,n}  \qquad \text{a.s.}
\end{equation}
Moreover, we clearly have from \eqref{DEFbr} that
\begin{equation*}
b_{r,i+1}= \left(\frac{r-i}{\alpha -i}\right)b_{r,i} \hspace{1cm} \text{and} \hspace{1cm} b_{r,i+1}-b_{r,i}= \left(\frac{r-\alpha}{\alpha -i}\right)b_{r,i}
\end{equation*}
Therefore, we deduce from \eqref{DEFPN}, \eqref{DEFIPr}, \eqref{DEFSNr}, \eqref{eq:esr1} together with straighforward calculation that
\begin{equation}
\label{eq:esr2}
\E[S_{r,n+1}\,|\,\mathcal{F}_{r,n}]=\gamma_{r,n}S_{r,n}-(1-\gamma_{r,n})\frac{(-1)^{r}p_{\alpha}(r) \theta}{\alpha}  \qquad \text{a.s.}
\end{equation}
where
\begin{displaymath}
\gamma_{r,n}=1-\left(\frac{r-\alpha}{n+\theta}\right).
\end{displaymath}
Hereafter, we obtain from \eqref{DEFMr} and \eqref{eq:esr2} and the elementary identity $a_{r,n+1}\gamma_{r,n}= a_{r,n}$
that for all $n\geq1$,
\begin{align*}
\E[M_{r,n+1}\,|\,\mathcal{F}_{r,n}]&=\E\left[a_{r,n+1}\left(S_{r,n+1}+\frac{(-1)^{r} p_{\alpha}(r) \theta}{\alpha}\right)\,\Bigr|\,\mathcal{F}_{r,n}\right] \qquad \text{a.s.}\\
&=a_{r,n+1}\left(\E[S_{r,n+1}\,|\,\mathcal{F}_{r,n}]+\frac{(-1)^{r} p_{\alpha}(r) \theta}{\alpha}\right) \qquad \text{a.s.}\\
&=a_{r,n+1}\left(\gamma_{r,n}S_{r,n}-(1-\gamma_{r,n})\frac{(-1)^{r}p_{\alpha}(r) \theta}{\alpha}+\frac{(-1)^{r} p_{\alpha}(r) \theta}{\alpha}\right) \qquad \text{a.s.}\\
&=a_{r,n+1}\gamma_{r,n} \left(S_{r,n}+\frac{(-1)^{r} p_{\alpha}(r) \theta}{\alpha}\right)  \qquad \text{a.s.}\\
&=a_{r,n} \left(S_{r,n}+\frac{(-1)^{r} p_{\alpha}(r) \theta}{\alpha}\right)= M_{r,n} \qquad \text{a.s.}
\end{align*}
Moreover, the sequence $(M_{r,n})$ is square integrable as for all $1 \leq i \leq r$, $K_{i,n}\leq K_n \leq n$. Consequently, $(M_{r,n})$ is a locally square
integrable martingale. Its predictable quadratic variation $\langle M_r \rangle_{n}$ is given by
\begin{equation}
\label{eq:CVMr1}
\langle M_r \rangle_{n} = \sum_{k=1}^{n-1} \E[(\Delta  M_{r,k+1})^{2}\,|\,\mathcal{F}_{r,k}]
\end{equation}
where $\Delta M_{r,n+1}=M_{r,n+1}-M_{r,n}$. It follows from \eqref{DEFAr}, \eqref{DEFMr} and \eqref{eq:esr2} that
\begin{equation}
\label{DECOMG}
\Delta M_{r,n+1}= a_{r,n+1} \Big( S_{r,n+1} - \E[S_{r,n+1}\,|\,\mathcal{F}_{r,n}]\Big)=a_{r,n+1} \Big( \zeta_{r,n+1} - \E[\zeta_{r,n+1}\,|\,\mathcal{F}_{r,n}]\Big)
\end{equation}
where
\begin{equation*}
\zeta_{r,n+1}= \sum_{i=1}^{r}b_{r,i}\xi_{i,n+1}+(-1)^r p_{\alpha}(r) \xi_{n+1}.
\end{equation*}
Consequently, we immediately obtain from decomposition \eqref{DECOMG} that
\begin{equation}
\label{eq:CVMr2}
\E[(\Delta M_{r,n+1})^{2}\,|\,\mathcal{F}_{r,n}]=a_{r,n+1}^{2}\Big(\E[\zeta_{r,n+1}^{2}\,|\,\mathcal{F}_{r,n}]-\E^2[\zeta_{r,n+1}\,|\,\mathcal{F}_{r,n}]\Big) \qquad \text{a.s.}
\end{equation}
We already saw from \eqref{eq:esr0} that
\begin{equation}
\label{CEzeta1}
\E[\zeta_{r,n+1}\,|\,\mathcal{F}_{r,n}]= \sum_{i=1}^{r}b_{r,i}(p_{i,n}-q_{i,n}) + (-1)^rp_{\alpha}(r) p_{n} \qquad \text{a.s.}
\end{equation}
In addition, we also have $\E[\xi_{n+1}^2\,|\,\mathcal{F}_{r,n}]=p_n$ and $\E[\xi_{i,n+1}^2\,|\,\mathcal{F}_{r,n}]=p_{i,n}+q_{i,n}$ for $1 \leq i \leq r$.
Moreover, $\E[\xi_{i,n+1} \xi_{j,n+1}\,|\,\mathcal{F}_{r,n}]=0$ as soon as $|i-j| \geq 2$. Furthermore, one can observe that $\E[\xi_{i,n+1} \xi_{i+1,n+1}\,|\,\mathcal{F}_{r,n}]=-q_{i, n}$ for $1 \leq i \leq r-1$ and
$\E[\xi_{i,n+1} \xi_{i-1,n+1}\,|\,\mathcal{F}_{r,n}]=-q_{i-1,n}$ for $2 \leq i \leq r$ and
$\E[\xi_{n+1} \xi_{i,n+1} \,|\,\mathcal{F}_{r,n}]=p_n$ for $i=1$, $\E[\xi_{n+1} \xi_{i,n+1} \,|\,\mathcal{F}_{r,n}]=0$ for $2 \leq i \leq r$. Hence,
we have 
\begin{equation}
\label{CEzeta2}
\E[\zeta_{r,n+1}^2\,|\,\mathcal{F}_{r,n}]= \sum_{i=1}^{r}b_{r,i}^2(p_{i,n}+q_{i,n}) -2\sum_{i=1}^{r-1}b_{r,i}b_{r,i+1}q_{i,n} - \Bigl(\frac{2r-\alpha}{\alpha}\Bigr) p_{\alpha}^2(r)p_n
\quad \text{a.s.}
\end{equation}
We have from the almost sure convergences \eqref{eq:kalpha} and \eqref{eq:kralpha} together with \eqref{DEFPN} and \eqref{DEFIPr} that
\begin{equation}
\label{ASCVGPN}
\lim_{n \rightarrow \infty} n^{1-\alpha} p_n= \alpha S_{\alpha, \theta} \hspace{1cm} \text{a.s.}
\end{equation}
and for all $i \geq 1$,
\begin{equation}
\label{ASCVGQIN}
\lim_{n \rightarrow \infty} n^{1-\alpha} q_{i,n}= (i-\alpha) p_{\alpha}(i) S_{\alpha, \theta} \hspace{1cm} \text{a.s.}
\end{equation}
However, for all $i \geq 2$, $p_{i,n}=q_{i-1,n}$. Hence, we obtain from \eqref{ASCVGQIN} that for all $i \geq 2$, 
\begin{align}
\label{ASCVGPQINDif}
\lim_{n \rightarrow \infty} n^{1-\alpha} (p_{i,n}- q_{i,n})
&= ((i-1-\alpha) p_{\alpha}(i-1) -(i-\alpha) p_{\alpha}(i)) S_{\alpha, \theta} \notag
\hspace{1cm} \text{a.s.} \\
&=\alpha p_{\alpha}(i) S_{\alpha, \theta} \hspace{1cm} \text{a.s.}
\end{align}
Therefore, it follows from \eqref{CEzeta1}, \eqref{ASCVGPN} and \eqref{ASCVGPQINDif} that
\begin{equation}
\label{ASCVGPCEzetar1}
\lim_{n \rightarrow \infty} n^{1-\alpha} \E[\zeta_{r,n+1}\,|\,\mathcal{F}_{r,n}]= \alpha  \Big( \sum_{i=1}^{r}b_{r,i} p_{\alpha}(i) + (-1)^r  p_{\alpha}(r) \Big)S_{\alpha, \theta} =0 \hspace{1cm} \text{a.s.}
\end{equation}
since by the Binomial theorem, the sum in \eqref{ASCVGPCEzetar1} reduces to
\begin{align}
\sum_{i=1}^{r}b_{r,i} p_{\alpha}(i) &= \sum_{i=1}^{r}
(-1)^{r-i}\frac{(i-\alpha)^{(r-i)}}{(r-i)!} \frac{\alpha(1-\alpha)^{(i-1)}}{i!} \notag \\
&=\frac{\alpha(1-\alpha)^{(r-1)}}{r!} \sum_{i=1}^{r} \binom{r}{i} (-1)^{r-i} \notag \\
&=-(-1)^r  p_{\alpha}(r).
\label{BinomialBP}
\end{align}
By the same token, we deduce from \eqref{CEzeta2}, \eqref{ASCVGPN} and \eqref{ASCVGQIN} that 
\begin{align}
\label{ASCVGPCEzetar2}
\lim_{n \rightarrow \infty} n^{1-\alpha} \E[\zeta_{r,n+1}^2\,|\,\mathcal{F}_{r,n}]
%&= \Big( \sum_{i=1}^{r}b_{r,i}^2 (2i - \alpha) p_{\alpha}(i) -2 \sum_{i=1}^{r-1}b_{r,i}b_{r,i+1}(i-%%\alpha)p_{\alpha}(i) - (2r-\alpha) p_{\alpha}^2(r) \Big)S_{\alpha, \theta} \notag \\
&= \Big( \sum_{i=1}^{r}b_{r,i}^2 ((2i - \alpha)+2(r-i)) p_{\alpha}(i) - (2r-\alpha)p_{\alpha}^2(r) \Big)S_{\alpha, \theta} \notag \\  
&= (2r- \alpha) \Big( \sum_{i=1}^{r}b_{r,i}^2  p_{\alpha}(i) - p_{\alpha}^2(r) \Big)S_{\alpha, \theta} 
\qquad \text{a.s.} \notag \\  
&=  p_{\alpha}(r) 
\frac{(r-\alpha)^{(r+1)}}{r!} S_{\alpha, \theta}  \qquad \text{a.s.}
\end{align}
since the Chu-Vandermonde identity for the rising factorial
\begin{equation}
\label{ChuVan}
\sum_{k=0}^n (-n)^{(k)} \frac{(b)^{(k)}}{(a)^{(k)} k!}
= \sum_{k=0}^n (-1)^{k} \binom{n}{k} \frac{(b)^{(k)}}{(a)^{(k)}}= \frac{(a-b)^{(n)}}{(a)^{(n)}}
\end{equation}
which holds for all $a,b \in \dR$ with $a \notin \{0,-1,-2,\ldots\}$, implies via \eqref{DEFbr} that
\begin{align*}
\sum_{i=1}^{r}b_{r,i}^2  p_{\alpha}(i)&=\sum_{i=1}^{r}\left(\frac{(i-\alpha)^{(r-i)}}{(r-i)!}\right)^{2}\frac{\alpha(1-\alpha)^{(i-1)}}{i!}\\
&=\Gamma^2(r-\alpha)\sum_{i=1}^{r}\frac{\alpha(1-\alpha)^{(i-1)}}{\Gamma^2(i-\alpha)((r-i)!)^{2}i!}\\
&=-\frac{\Gamma^2(r-\alpha)}{\Gamma^2(-\alpha)(r!)^{2}}\sum_{i=1}^{r}\frac{(-r)^{(i)}(-r)^{(i)}}{(-\alpha)^{(i)}i!}\\
&=-p_\alpha^2(r)\left(\frac{(r-\alpha)^{(r)}}{(-\alpha)^{(r)}}-1\right).
\end{align*}
Hereafter, we obtain from standard results on the asymptotic behavior of
the Euler Gamma function together with \eqref{DEFAr} that
\begin{equation}
\label{CVGANr}
\lim_{n\rightarrow+\infty} \frac{a_{r,n}}{n^{r-\alpha}}=\lim_{n\rightarrow+\infty}\frac{1}{n^{r-\alpha}}
\frac{(\theta+1)^{(n-1)}}{(\alpha+\theta-r+1)^{(n-1)}}=\frac{\Gamma(\alpha+\theta-r+1)}{\Gamma(\theta+1)}.
\end{equation}
Therefore, it follows from \eqref{eq:CVMr2}, \eqref{ASCVGPCEzetar1}, \eqref{ASCVGPCEzetar2} and \eqref{CVGANr} that
\begin{equation*}
\lim_{n\rightarrow+\infty} \frac{\E[(\Delta M_{r,n+1})^{2}\,|\,(\mathcal{F}_{r,n})]}{n^{2r-\alpha -1}}=
p_{\alpha}(r) \frac{(r-\alpha)^{(r+1)}}{r!} \left(\frac{\Gamma(\alpha+\theta-r+1)}{\Gamma(\theta+1)}\right)^2 S_{\alpha, \theta}  \qquad \text{a.s.}
\end{equation*}
which implies through a direct application of Toeplitz lemma that
\begin{equation*}
\lim_{n\rightarrow+\infty}\frac{\langle M_{r}\rangle_{n}}{n^{2r-\alpha}}=p_{\alpha}(r)\frac{(r-\alpha)^{(r)}}{r!}\left(\frac{\Gamma(\alpha+\theta-r+1)}{\Gamma(\theta+1)}\right)^{2}S_{\alpha,\theta}\qquad\text{a.s.}
\end{equation*}
completing the proof of Lemma \ref{L-KEY}.
\end{proof}

\end{document}
